\section{Глава 1. Введение}
\subsection{Упражнение 1.1. Игра с самим собой(с.38)}
Предположим, что вместо игры со случайным противником описанный выше алгоритм обучения с подкреплением играет против себя самого, и обе стороны обучаются. Как вы думаете, что произойдет в таком случае? Обучится ли игрок другой стратегии выбора ходов?
\begin{solution}
    Заметим, что в таком случае мы имеем дело с нестационарной средой, в которой вероятность перехода в состояние $s'$ из состояния $s$ зависит от политики противника. \\
    Обратимся к словам автора на странице 36: "Если размер шага не уменьшается со временем до нуля, то этот игрок будет хорошо играть и с противниками, которые медленно изменяют манеру игры."\ Это означает, что если потребовать определенные условия на коэффициент $\alpha$, то такой способ обучения приведет к более универсальной стратегии, но лишь в "практическом"\ смысле, а не строгом математическом. О сходимости к оптимальной стратегии тут речь не идет, потому что классические результаты TD-обучения требуют стационарности среды.
\end{solution}

\subsection{Упражнение 1.2. Симметрии(c.38)}
Многие позиции в крестиках-ноликах выглядят различными, но на самом деле отличаются только симметрией. Как можно изменить описанный выше процесс обучения, чтобы воспользоваться этим фактом? Как это изменение могло бы улучшить процесс обучения? А теперь подумайте еще раз. Предположим, что противник не пользуется преимуществами симметрии. Должны ли тогда это делать мы? И верно ли, что позиции, совпадающие с точностью до симметрии, имеют одну и ту же ценность? 
\begin{solution}
    Сначала рассмотрим случай, когда оба игрока придерживаются принципа симметрии. В таком случае логично будет сгруппировать положения доски относительно некоторых канонических состояний. И обновлять ценности следует одновременно по всем эквивалентным позициям. Это приводит к меньшему количеству уникальных состояний, более быстрому обучению, более стабильным оценкам $V(s)$. \\
    Если же противник не пользуется симметрией, то совпадающие позиции перестают иметь одну и ту же ценность, потому что два симметричных состояния могут привести к разным распределениям будущих ходов. И, следовательно, пользоваться симметрией нам теперь невыгодно, так как предложенный алгоритм основан на значении $V(s)$.
\end{solution}

\subsection{Упражнение 1.3. Жадная игра(с.39)}
Предположим, что обучающийся с подкреплением игрок жадный, т. е. всегда выбирает ход, который переводит его в позицию с наивысшим рейтингом. Мог бы он научиться играть лучше нежадного игрока? Или хуже? Какие проблемы могли бы возникнуть?
\begin{solution}
    Абсолютное отсутствие исследования - ключевой недостаток жадного подхода. Теоретически, агент может научиться играть лучше нежадного игрока, но по чистой случайности, в искусственных случаях. В практических задачах это никогда не встретится. Исключение возможности исследовать пространство решений приводит к застреванию в локальных областях, далеко не оптимальным стратегиям. Объясняется это тем, что игрок навсегда закрепляет ошибочные ранние оценки. Ещё одна проблема - зависимость от начальных условий.
\end{solution}

\subsection{Упражнение 1.4. Обучение на результатах разведки(с.39)}
Предположим, что обучение подвергается обновлению после любого хода, включая разведочные. Если со временем размер шага уменьшается подходящим образом (но не уменьшается стремление к исследованию), то ценности состояний будут сходиться к другому набору вероятностей. Что концептуально представляют собой два набора вероятностей, вычисленных, когда мы обучаемся и не обучаемся на разведочных ходах? В предположении, что мы продолжаем совершать разведочные ходы, какой набор вероятностей предпочтительнее обучить? Какой принесет больше выигрышей?
\begin{solution}
    Метод без обучения на разведочных данных вычисляет ценности оптимальной политики. Это идеальные ожидаемые выигрыши, если бы агент всегда действовал оптимально, без исследовательских шагов. Но при этом агент совершает случайные действия. Это означает, что найденные ценности соответствуют не той политике, которой придерживается игрок, а идеальному случаю. \\
    Обновление ценностей на разведочных ходах соответствует реальному поведению агента, поэтому логичнее их учитывать, чтобы вычисленные ценности соответствовали действительности. При достаточном обучении именно такой подход принесет больше выигрышей по этой же причине, несмотря на выполнение следующего неравенства:
    \[v_*(s) \ge v_\pi(s) \qquad \forall s \in S,\]
    где $v_*(s)$ - ценность состояния $s$ при оптимальной политике, а $v_\pi(s)$ - ценность состояния $s$ при политике, включающей разведочные ходы.
\end{solution}

\subsection{Упражнение 1.5. Другие улучшения(с.39)}
Можете ли вы предложить иные способы улучшить игрока, обучающегося с подкреплением? Можете ли вы придумать лучший способ решения задачи об игре в крестики-нолики в том виде, в котором она поставлена?
\begin{solution}
   Кроме улучшений, предложенных в предыдущих упражнениях, я могу предложить следующие:
   \begin{itemize}
       \item Обновление не только предыдущего шага, но и всей истории по правилу
        \[
            V(s_{t-k}) \leftarrow V(s_{t-k}) 
            + \alpha_k \bigl( V(s_{t+1}) - V(s_{t-k}) \bigr),
            \quad k=\overline{0,t} \quad\alpha_k \downarrow \text{по } k
        \]
       \item Постепенное уменьшение вероятности исследования с течением времени улучшает итоговую игру.
   \end{itemize}
   Крестики-нолики являются антагонистической игрой(выигрыш одного равен проигрышу другого) с полной информацией. Более того, множество состояний конечно и невелико. Поэтому к игре применимы классические подходы из теории игр, исследующие минимаксные стратегии.
\end{solution}
